\section{Future work: optimising the soft-HMC controller}
\label{sec:future-work}

The chance-constrained variants of Section~\ref{sec:control-cost-variants}
settled the question of \emph{which} cost surrogate to use: the soft
sigmoid surrogate driven by HMC. This section is about the question
that remains: how to make the soft-HMC controller faster without
sacrificing the control quality that justified the choice in the
first place. The hard-IS variant is not in scope here --- it is a
documented dead-end (Section~\ref{sec:control-cost-variants}, the
``empirical comparison'' remark) and is preserved only for regression
testing.

\subsection{The trade-off as observed today} \label{sec:fw-tradeoff}

Running the soft-HMC controller's tempered SMC$^2$ with high-quality
settings --- more leapfrog steps per HMC trajectory, more HMC moves
per tempering level, finer tempering schedule --- produces better
control plans but costs more wall-clock time. This is, currently, the
framework's main practical bottleneck. Planning latency scales
roughly linearly with HMC work per tempering level:
\[
\text{(wall-clock per plan)}
\;\sim\; n_\lambda \cdot n_{\rm mcmc} \cdot n_{\rm leap}
\cdot \text{(grad eval cost)},
\]
where $n_\lambda$ is the number of tempering levels, $n_{\rm mcmc}$ is
the number of HMC moves per level, and $n_{\rm leap}$ is the leapfrog-
step count per move. All three factors trade quality against speed,
and the right operating point depends on the application's planning-
latency budget. The remaining subsections collect the levers and the
open directions for moving the Pareto frontier.

\subsection{The tunable knobs and what each one costs} \label{sec:fw-knobs}

Four configuration fields are first-line levers. Their definitions and
default values are in the configuration reference (Section~\ref{sec:config});
this subsection is about the quality--speed direction of each.

\begin{center}\small
\begin{longtable}{lp{3.5cm}p{3.5cm}p{4cm}}
\toprule
Knob & Quality direction & Speed direction & Diminishing-return knee \\
\midrule
\endhead
\texttt{hmc\_num\_leapfrog} & Higher = better mixing per trajectory; the proposal travels further per move. & Lower = fewer gradient evaluations per move. & Acceptance plateau around the standard $\tau \approx \pi/(2\omega)$ leapfrog tuning; raising past that wastes compute. \\
\texttt{num\_mcmc\_steps}     & Higher = more rejuvenation per tempering level, lower autocorrelation between tempering levels. & Lower = fewer rejuvenation moves per level. & Per-level autocorrelation typically halves with each doubling up to about $5$--$8$, then flattens. \\
\texttt{target\_ess\_frac}    & Higher (closer to $1$) = smaller $\delta\lambda$ per step, gentler tempering, more reliable convergence. & Lower = larger $\delta\lambda$, fewer tempering levels. & Below $\sim 0.3$ the bisection routinely fails to find a feasible step and the run aborts. \\
\texttt{max\_lambda\_inc}     & Smaller = caps each tempering jump tighter, more steps but each one safer. & Larger = bisection allowed to take bigger leaps. & Larger than the empirical first-window mean of $\delta\lambda$ has no effect; smaller than $\sim 0.01$ produces redundant levels. \\
\bottomrule
\end{longtable}
\end{center}

The four knobs interact. Lowering \texttt{target\_ess\_frac} to cut
the tempering-level count typically requires raising
\texttt{num\_mcmc\_steps} to compensate for the larger $\delta\lambda$
per step --- otherwise the cloud cannot keep up with the temperature
change and HMC acceptance collapses. The combinations that survive
this trade are application-specific.

\subsection{Open directions} \label{sec:fw-directions}

The following are candidate avenues for improving the soft-HMC
controller's Pareto frontier. None has been benchmarked; they are
listed in roughly decreasing order of how cheaply they could be
explored.

\paragraph{Profile first.} Before changing any algorithmic knob, run
\texttt{jax.profiler.start\_trace} around a representative call to
\path{smc2fc/control/tempered_smc_loop.py} and view the trace in
TensorBoard. The wall-clock split between (i) HMC leapfrog gradient
evaluations, (ii) the inner-PF likelihood, (iii) the on-device ESS
bisection, and (iv) host--device transfer latency is the precondition
for any informed knob change. The cheapest improvement may simply be
removing a host sync that nobody noticed.

\paragraph{Adaptive HMC step-size schedule along $\lambda$.} The
existing diagonal-mass-matrix re-estimation
(Section~\ref{sec:mass-matrix}) tightens the geometry per tempering
level, but the step size $\epsilon$ stays fixed. A schedule that
starts with a larger $\epsilon$ at low $\lambda$ (where the target is
nearly the prior, broad and forgiving) and shrinks toward a target
acceptance rate of $\approx 0.7$ at high $\lambda$ (where the
posterior is tight and the geometry is most curved) could reduce the
required $n_{\rm leap}$ at the easy end without sacrificing
acceptance at the hard end. BlackJAX has a dual-averaging adapter
that targets a fixed acceptance probability; wiring it into the
JAX-native loop is straightforward.

\paragraph{Per-tempering-level adaptive leapfrog count.} Currently
\texttt{hmc\_num\_leapfrog} is a single integer applied at every
tempering level. The optimal trajectory length depends on the local
geometry and on how far apart consecutive tempering levels are; a
schedule that tracks the per-level $\delta\lambda$ chosen by the
bisection (longer trajectories when $\delta\lambda$ is large) would
let the controller spend its leapfrog budget where it matters.

\paragraph{Reuse the inference posterior as the controller's prior.}
The controller's tempered SMC$^2$ currently uses a fixed Gaussian
prior $\thetactrl \sim \Ncal(0, \sigma^2 I)$ regardless of what the
filter has learned about $\thetadyn$. When the filter posterior is
informative --- in particular, when its uncertainty has collapsed
along the directions the controller cares about --- starting from a
prior that already encodes that information would shorten the
tempering schedule. The technical question is: what is the right
mapping from the filter's $\thetadyn$ posterior to the controller's
$\thetactrl$ prior? A heuristic such as ``set
$\sigma_{\rm prior}^2$ proportional to the filter's posterior
uncertainty along each direction'' is a starting point; a principled
mapping would require some thought about the joint problem.

\paragraph{Per-bridge HMC budget tuning.}
\texttt{num\_mcmc\_steps\_bridge} is currently $3$ (vs $5$ for cold
start). The bridge initialisation
(Section~\ref{sec:bridge}) is by construction close to the new
posterior, so fewer rejuvenation moves should be sufficient. The open
question is whether $3$ is the right number or whether $1$--$2$ would
do, and whether the answer depends on the magnitude of the
window-to-window posterior shift. A simple sweep on a representative
multi-window benchmark would settle it.

\paragraph{Mixed-precision HMC.} The HMC inner loop is dominated by
matrix--vector products in the gradient evaluation; the outer SMC$^2$
ESS arithmetic and tempering bisection are not. Running HMC in
\texttt{float32} and the rest of the pipeline in \texttt{float64} is
a common $\sim 2\times$ speedup on GPU. The risk is that for
ill-conditioned posterior geometries the lower precision degrades
HMC acceptance enough to require more moves per level, eating the
gain. Worth benchmarking.

\paragraph{Smaller controller particle count with re-warm-starts.}
The controller currently uses the same \texttt{n\_smc\_particles} as
the inference side. The control posterior is typically lower-dimensional
($d_c \approx 8$ vs $d_\theta \approx 16$ for typical applications)
and arguably needs fewer particles to characterise its mean
adequately. A reduction to $128$ or $96$ controller particles, paired
with the next window's bridge warm-start (Section~\ref{sec:bridge}
applied to the controller, not just to the filter), would scale the
per-window controller cost down proportionally. This requires
extending the bridge mechanism to the control posterior, which the
current code does not do.

\subsection{What's not on this list}

A few things deliberately not in the future-work list:

\begin{itemize}
  \item Replacing HMC with NUTS. Documented dead-end; see
        Section~\ref{sec:nuts-failed}.
  \item Replacing the soft surrogate with the hard indicator.
        Documented dead-end; see Section~\ref{sec:control-cost-variants}.
  \item Replacing the Gaussian bridge with the SF/BW variant by
        default. Documented as not paying for itself; see
        Section~\ref{sec:bridge}.
\end{itemize}

These are not open questions. Future contributors should consult the
relevant section before re-attempting.
